\paragraph{A trial quantum history with a global norm bound.}
The complete reduced metric also permits an explicit, error-controlled
trial history in the full interacting Hilbert space. The bound below is
deliberately conservative: its time interval is extremely short in the
supplied dimensionless units. It does not furnish an ordinary-physics
benchmark, a physical clock, or an exact computed Hamiltonian trajectory.

Retain the same fixed cone and set
\(e=1/4\), \(m^2=1/2\), \(g=1/4\), \(\hbar=1\), and
\(\sigma=1/2\). Let \(H\) be the neutral Hamiltonian of
Theorem~\ref{thm:whitney-interacting-quantum} and \(f=f_\sigma\) the
normalized state of Proposition~\ref{prop:whitney-interacting-gaussian-state}.
All configuration coordinates remain present:
\(q=(a,\psi)\in\mathcal Q\cong\mathbb R^{30}\times\mathbb C^{13}\),
with Euclidean-orthonormal coordinates on the Coulomb subspace. Write
\(X=|a|^2\), \(Y=|\psi|^2\), and \(\rho=\sqrt{\det\gamma}\).

\begin{theorem}[Controlled potential-phase trial]
\label{thm:whitney-quantum-trial-history}
The normalized neutral trial path
\[
 v(t,q)=e^{-itV(q)/\hbar}f(q)
\]
belongs to \(\mathcal D(H)\) and is differentiable in the Hilbert norm.
If \(u(t)=e^{-itH/\hbar}f\), then
\begin{equation}
 \|u(t)-v(t)\|_{L^2(d\mathrm{vol}_\gamma)}^2
 \leq\frac{(\overline E+\overline K)|t|
                   +\overline B|t|^3/6}{\hbar},
 \label{eq:whitney-quantum-trial-error}
\end{equation}
where the following rational constants are admissible:
\begin{align*}
 \overline K&=135829605095437,&
 \overline E&=\frac{2173273682146345}{16},&
 \overline B&=\frac{24441816695556087}{4}.
\end{align*}
In particular, \(\|u(t)-v(t)\|\leq1/10\) for every
\(|t|\leq2^{-55}\). The trial configuration probability measure is
exactly \(N(0,\sigma^2I_{56})\) at every time; its phase changes, while
its configuration density stays fixed.
\end{theorem}

\begin{proof}
The argument separates a form comparison from explicit coefficient bounds.
Define
\[
 K_0=\frac{\hbar^2}{2}\int|df|_{\gamma^{-1}}^2\,d\mu,
 \qquad E_0=K_0+\int V|f|^2\,d\mu,
 \qquad B_0=\int|dV|_{\gamma^{-1}}^2|f|^2\,d\mu.
\]
The polynomial derivative bounds used for the initial state also apply to
\(e^{-itV/\hbar}f\) on bounded time intervals. Thus \(v(t)\), its
formal Hamiltonian image, and \(Vv(t)\) lie in \(L^2(d\mu)\).
The adjoint-domain characterization and essential self-adjointness put
\(v(t)\) in \(\mathcal D(H)\); differentiation under its Gaussian
envelope gives \(i\hbar\dot v=Vv\). Gauge invariance of \(V\) gives
neutrality, and the phase has unit modulus.

Let \(\mathfrak t\) denote the kinetic quadratic form. Since \(f\) is
real and positive,
\(\mathfrak t(v(s),v(s))=K_0+s^2B_0/2\). Positivity of \(V\)
and conservation of the exact Hamiltonian energy give
\(\mathfrak t(u(s),u(s))\leq E_0\). Differentiating the overlap and
using the form Cauchy inequality yields
\[
 \left|\frac{d}{ds}\|u(s)-v(s)\|^2\right|
 \leq\frac{2}{\hbar}\sqrt{E_0}\sqrt{K_0+s^2B_0/2}
 \leq\frac{E_0+K_0+s^2B_0/2}{\hbar}.
\]
Integration proves \eqref{eq:whitney-quantum-trial-error} for any valid
upper bounds on these three constants. This comparison uses the full
kinetic form and its Riemannian measure.

The canonical geometry has volume at most \(18\), tetrahedron volumes
at least \(5/6\), squared edge lengths at most \(4\), and
\(|\nabla\lambda_i|\leq2\). Every node belongs to at least five
tetrahedra and every edge to at least two. These inequalities are checked
in \(\mathbb Q(\sqrt5)\) from the canonical vertices and incidence lists.
They imply the useful uniform bounds
\begin{equation}
 M\geq\frac1{81}I,\qquad W(a)^*W(a)\geq\frac1{128}I.
 \label{eq:whitney-quantum-history-coercivity}
\end{equation}
For the first, an affine vector field \(F\) on a tetrahedron satisfies
\(\int_T|F|^2\geq |T|\sum_i|F(v_i)|^2/20\), while its six edge
integrals obey \(\sum_{ij}|a_{ij}|^2\leq6\sum_i|F(v_i)|^2\).
Summing gives \(M\geq1/72\), which implies the displayed bound.
For the second, choose a largest local nodal coefficient \(|z_i|\).
On \(\lambda_i\geq1/2\),
\(|Wz|\geq(2\lambda_i-1)|z_i|\), independently of the phases, and
\[
 \frac1{|T|}\int_{\lambda_i\geq1/2}(2\lambda_i-1)^2\,dx
 =3\int_{1/2}^1(2s-1)^2(1-s)^2\,ds=\frac1{80}.
\]
Thus the summed lower bound is \(5/384\), which exceeds \(1/128\).

The radial edges give \(|D\xi|^2\geq|\xi|^2\) for mean-zero
\(\xi\): their contribution is
\(\sum_{i=1}^{12}(\xi_i-\xi_0)^2=|\xi|^2+13\xi_0^2\).
The mean-zero vertical Gram matrix \(I\) consequently satisfies
\(I\geq1/81\). For a slice velocity \((b,z)\) and mean-zero
gauge velocity \(\eta\), let
\(E_\eta=G(w+R\eta,w+R\eta)\). Maxwell orthogonality gives
\(|b|^2+|\eta|^2\leq81E_\eta\). The vertical interpolation identity
and \(\|W\|\leq\sqrt{18}\), \(\|S\|\leq\sqrt{18}|\psi|/4\)
give
\(\|Wz\|^2\leq(1+729Y/2)E_\eta\). Applying
\eqref{eq:whitney-quantum-history-coercivity} and minimizing over
\(\eta\) proves
\begin{equation}
 \gamma^{-1}\leq A(Y)I,\qquad A(Y)=209+46656Y.
 \label{eq:whitney-quantum-history-inverse}
\end{equation}

The density derivative can be bounded without differentiating a matrix
inverse twice. The combined slice and vertical tangent matrix has constant
determinant: after permuting columns its edge block is the fixed matrix
\([T_c,DB]\) and its scalar diagonal block is the identity. Schur
factorization therefore gives
\(\log\rho=\mathrm{constant}+\tfrac12\log\det G-\tfrac12\log\det I\).
Put \(J=(S,W)\). For any full velocity \(w=(w_a,w_\psi)\),
\eqref{eq:whitney-quantum-history-coercivity} implies
\[
 |w_a|\leq9\sqrt{G(w,w)},\qquad
 |w_\psi|\leq(8+108|\psi|)\sqrt{G(w,w)}.
\]
Each real phase-path row has Euclidean norm at most one. Differentiating
\(J\) and using \(2\|Jw\|^2\leq G(w,w)\) yields, for a
configuration direction \(h\),
\[
 |dG[h](w,w)|\leq3(17+441|\psi|/4)|h|G(w,w).
\]
For \(\chi=W_0B\), the identity
\(I=I_0+2e^2\int\chi^{\mathsf T}\chi|\Psi_a|^2\)
similarly gives
\(|dI[h](\eta,\eta)|\leq27(1+|\psi|/4)|h|I(\eta,\eta)\).
Taking traces in dimensions \(68\) and \(12\) proves
\begin{equation}
 |\nabla\log\rho|\leq1896+11286\sqrt Y.
 \label{eq:whitney-quantum-history-density}
\end{equation}
All gradients here are in the orthonormal slice coordinates; bounds for
unit full-configuration directions restrict to those coordinates.

For completeness, the remaining potential bounds follow directly from
the same geometry. Whitney curl has squared operator bound \(384\)
on each cell. The pointwise estimates
\begin{align*}
 |\Psi_a|&\leq|\psi|,&
 |D_x\Psi_a|&\leq(4+9|a|/4)|\psi|,\\
 |d(D_x\Psi_a)[b,z]|&\leq(4+9|a|/4)|z|
                   +(13/4+9|a|/16)|\psi||b|
\end{align*}
retain the derivative of the gauge-dependent dressing. The phase-gradient
bound uses \(|\nabla p_i\cdot a|\leq4|a|\) and
\(|\mathbf A(a)|\leq5|a|\). Squaring and integrating gives
\begin{align*}
 V&\leq V_*(X,Y):=3456X+585Y+\frac{729}{4}XY+\frac94Y^2,\\
 |\nabla V|^2&\leq P(X,Y):=
 4\cdot6912^2X+2313^2Y^2+4\cdot1152^2X^2Y^2+\frac{81}{4}Y^4\\
 &\hspace{38mm}+3\cdot1170^2Y+3\cdot1152^2X^2Y+243Y^3.
\end{align*}
For example, the vector and scalar gradient norms are bounded by
\(6912\sqrt X+(2313/2)Y+1152XY+(9/4)Y^2\) and
\(1170\sqrt Y+1152X\sqrt Y+9Y^{3/2}\); Cauchy's inequality for
four and three summands gives the displayed \(P\).

Under \(|f|^2d\mu\), the independent variables \(X,Y\) have laws
\(\sigma^2\chi^2_{30}\) and \(\sigma^2\chi^2_{26}\). The Gaussian
integration-by-parts recurrence gives exactly
\[
 \mathbb E[(\sigma^2\chi^2_d)^k]
       =\sigma^{2k}\prod_{j=0}^{k-1}(d+2j).
\]
Since \(df=-f(d\log\rho+q\cdot dq/\sigma^2)/2\), admissible
constants are the finite polynomial expectations
\begin{align*}
 \overline K&=\frac14\mathbb E\!\left[
 A(Y)\{2\cdot1896^2+2\cdot11286^2Y+16(X+Y)\}\right],\\
 \overline E&=\overline K+\mathbb E[V_*]
             =\overline K+\frac{619353}{16},&
 \overline B&=\mathbb E[A(Y)P(X,Y)].
\end{align*}
Exact rational evaluation gives the stated values. The right side of
\eqref{eq:whitney-quantum-trial-error} is increasing in \(|t|\);
at \(2^{-55}\) it is less than \(0.007541<1/100\), whereas the
preceding dyadic time \(2^{-54}\) fails that target. This proves coverage
of the entire stated interval, including negative times.
\end{proof}

The packet
\path{code/electromagnetism/runtime/whitney_quantum_history_receipt.json}
records five trial times and four configuration probes. Potential values
and the relative phase angles at those probes are exact elements of
\(\mathbb Q(\sqrt5)\); an independent verifier reconstructs their
simplex integrals and all envelope constants. The sampled phases illustrate
the trial function; the norm bound concerns the full Hilbert space.
Global polynomial integration includes the entire Gaussian tail, with no
Monte Carlo, quadrature, or floating-point allowance used to certify the
error. A separate numerical evaluator retains the full metric density and
marks its approximate coefficients explicitly. The certificate inherits the
supplied action, geometry, quantum prescription, initial state and time
units. Its tiny conservative horizon is not physical state preparation,
authenticated observer activity, spatial convergence, or an interacting
continuum quantum field theory.
